% !TeX root = ../thuthesis-example.tex

\chapter{引言}

\section{项目背景}

近年来，随着移动互联网基础设施的持续完善、冷链物流技术的稳步升级以及消费者购物习惯的深刻转变，生鲜电商行业经历了前所未有的高速增长。据行业统计数据显示，2024年中国生鲜电商市场规模已突破6000亿元人民币，年均复合增长率保持在20\%以上，其中社区团购、即时零售与前置仓模式成为推动行业扩张的三大主要驱动力。以线上订购、即时配送为代表的生鲜水果零售场景，因其高频消费、强时效性和对用户体验的极致追求，已成为电商领域中技术挑战最为密集的垂直赛道之一。生鲜系统不仅需要支持用户在移动端的高频浏览、快速下单、购物车管理与即时支付，还需要保障后台运营人员对商品库存、分类管理、订单流转及多维度运营数据的实时掌控。

然而，传统的单体架构在面对上述业务场景时暴露出了一系列结构性缺陷。在部署效率层面，单体应用将所有业务模块打包为一个整体制品，任何一个模块的微小变更都需要对整个系统重新编译与发布，部署周期往往长达数十分钟，严重制约了业务的快速迭代。在弹性扩容层面，单体架构无法针对特定高负载模块进行独立的水平扩展，促销高峰期只能整体扩容，造成计算资源的严重浪费。在数据库层面，单体架构通常采用单一数据库实例承载全部业务数据，随着并发查询量的增长，数据库极易成为系统的单点瓶颈，表现为查询延迟攀升与锁竞争加剧。在故障隔离方面，各模块共享运行时环境，任何一个模块的内存泄漏或未捕获异常都可能导致整个系统崩溃。在团队协同层面，多人在同一代码仓库中修改不同模块时，代码冲突频繁发生，模块间的隐式依赖也使得回归测试范围难以界定。

分布式微服务架构为上述问题提供了系统性的解决思路。根据微服务架构设计原则，通过将庞大的业务应用按照功能边界拆分为多个独立部署的微服务单元，配合统一网关路由、服务发现注册中心、声明式服务调用、分布式缓存、消息队列异步解耦、对象存储及长连接实时推送，能够显著提升系统的弹性扩容能力与开发迭代效率。为此，本项目以"生鲜"生鲜水果运营系统为业务蓝本，在已有的Spring Boot核心业务逻辑基础上，开展了课程设计级别的轻量微服务化改造与分布式应用治理。系统既能够保证本地单机运行环境下的极高演示稳定性，又全面引入了Spring Cloud及相关分布式中间件的技术底座，使之完全契合分布式应用开发技术课程的实践与学术验收要求。本选题的核心意义在于，它为团队成员提供了将课堂分布式理论知识转化为工程实践能力的完整平台，通过真实业务场景驱动，使每位成员深入理解微服务拆分策略、服务治理机制与工程化质量保障体系之间的内在关联，为后续的毕业设计和企业级项目开发积累了宝贵的实战经验。

\section{项目目标}

本项目旨在构建一个高可用、松耦合且具备完整工程化测试验证闭环的生鲜分布式运营系统，其总体目标可从核心业务闭环、微服务架构治理、分布式中间件集成以及工程化质量保障四个维度进行阐述。

在核心业务闭环目标方面，系统致力于实现从商品分类管理、品种维护、上下架状态控制，到小程序端用户浏览商品、加购操作、账单生成以及模拟微信支付的完整业务流转闭环。选择覆盖商品全生命周期管理与交易全链路闭环，是因为只有打通从供给侧的商品录入到需求侧的支付确认这一完整路径，才能真实验证分布式架构在实际业务场景中的数据一致性保障效果。预期管理员可在后台完成商品全流程管理，消费者可在小程序端体验从浏览到支付 of 无缝流程，两端数据通过微服务接口实现实时同步。

在微服务架构治理目标方面，系统采用Spring Cloud Nacos同时承担服务注册发现与动态配置管理的双重职责。选择Nacos是因为其原生支持配置管理功能，能够将注册发现与配置中心合二为一，减少中间件的部署复杂度。系统通过Spring Cloud Gateway构建统一网关路由入口，负责全平台流量的路径匹配转发与跨域请求拦截。选择Gateway是因为其基于WebFlux响应式编程模型，在高并发场景下具有更优的线程利用率，且断言与过滤器机制的可编程性更强。预期所有外部请求统一经由网关进入系统内部，各微服务通过Nacos实现动态发现与负载均衡，配置变更无需重启即可生效。

在分布式中间件集成目标方面，系统利用OpenFeign声明式调用实现跨服务的强类型数据同步，通过接口注解将远程HTTP调用抽象为本地方法调用，降低跨服务通信的编写复杂度。集成RabbitMQ消息队列将订单支付成功事件与后台通知推送解耦，采用RabbitMQ是因为其路由机制灵活，Exchange与Binding的组合模式能够精确控制消息投递路径，在中小规模消息吞吐场景下部署门槛更低。通过WebSocket全双工通道实现语音与卡片式即时播报提醒，使运营人员能够第一时间感知新订单的产生。预期核心链路通过同步Feign调用保障一致性，非关键通知链路通过异步消息队列解耦，实时推送通道保障信息即时触达。

在工程化质量保障目标方面，项目采用PM2进行长运行进程守护和端口资源锁定，排除演示现场服务崩溃隐患。基于Newman开展覆盖全量RESTful端点的接口自动化测试，确保每个API接口的参数校验、状态码断言与返回数据验证均纳入自动化回归范围。结合Playwright建立包含登录绕过、直接落库与消息响应检测的端到端质量验证体系，模拟真实用户操作路径进行全链路验证。保留Docker Compose容器化编排预案作为生产环境的后续演进通道。预期通过多层次质量保障手段，确保系统在各种运行环境下均能稳定完成业务演示与技术验收。

\section{团队分工与组织结构}

为了保证微服务接口的严谨交付，第五组的四位成员在项目生命周期中开展了紧密协同。在代码管理方面，团队采用Git进行版本控制，遵循feature分支开发策略，每位成员在功能分支上独立开发，完成后通过Pull Request发起合并，由组长进行代码审查后方可合入主分支。在沟通协调方面，团队在开发高峰期保持每日站会机制，及时发现并解决跨服务接口对接中的兼容性问题。在文档管理方面，团队统一使用仓库docs目录进行需求文档、设计文档与测试报告的集中管理。

组长曹磊承担系统架构师与项目总协调人角色。需求分析阶段主导完成功能需求梳理与微服务边界划分。设计阶段完成Gateway路由规则设计、Nacos命名空间规划及PM2编排方案设计。编码阶段独立完成lgg-gateway模块全部代码，包括8个路由断言规则、3个全局过滤器及跨域配置，搭建Nacos配置文件共计12份，编写PM2生态系统配置文件与多服务启动脚本。测试阶段统筹设计Playwright端到端自动化架构，编写5个核心测试用例。此外负责课程设计报告的统稿与排版工作。

陈佳明担任核心业务开发工程师，负责lgg-business模块的重构。需求分析阶段参与商品管理、分类管理与订单管理的功能细化与接口定义。编码阶段完成三大模块共计15个RESTful接口的持久层CRUD改造与业务逻辑实现，涵盖Controller层、Service层与Mapper层的全部编写。基于Spring AOP实现公共字段自动填充切面，通过自定义注解实现createTime、updateTime等审计字段的统一注入，涉及切面类1个、自定义注解1个及6处方法标记。测试阶段配合完成业务接口的Postman测试用例编写与调试。

陈晓楠独立负责lgg-pay支付微服务的设计与开发。需求分析阶段完成模拟支付场景的需求拆解与边界条件分析。编码阶段独立完成6个支付接口的全部代码，包括模拟下单、支付回调与订单状态查询等功能。针对Spring Security拦截链进行定制放行配置，确保支付回调接口能绕过认证拦截。完成OpenFeign远程调用客户端封装，定义Feign接口用于跨服务数据同步，并设置合理的超时机制保障容错能力。测试阶段独立完成支付模块全部接口的功能验证与异常场景测试。

程志阳专注于lgg-notice通知服务开发，承担通知开发工程师与测试工程师双重角色。编码阶段独立实现4个通知接口与2个WebSocket端点的全部代码，设计高可靠性的RabbitMQ异步消费队列，包括交换机声明、队列绑定、消息监听器与确认机制的完整实现，构建基于WebSocket的实时广播通道，支持语音播报与卡片式消息两种推送形态。测试阶段负责编写覆盖全部微服务核心接口的Newman接口自动化测试集，包含请求参数构造、响应断言脚本与测试数据管理的完整测试工程。

\section{本文组织结构}

本文后续内容按照技术基础、系统设计、功能实现、测试验收与总结展望的逻辑主线依次展开，各章之间形成由理论到实践、由整体到细节、由建设到验证的递进关系。第2章为技术基础与相关工作章节，系统介绍Spring Boot、Spring Cloud、Nacos、Gateway、OpenFeign、RabbitMQ、WebSocket、PM2、Newman及Playwright等核心技术的基本原理与在本系统中的定位，为后续章节提供理论支撑与技术语境。第3章为系统分析与设计章节，详述功能需求分析与非功能需求约束，完成微服务边界定义与服务拆分策略论证，给出系统架构图、服务交互时序图及数据库表结构设计方案，该章节是连接技术基础与功能实现的关键桥梁。第4章为系统实现章节，作为全文核心章节，重点展示管理端运营大屏、商品管理、支付闭环、WebSocket通知与PM2运维管理的具体编码实现，将设计方案转化为可运行代码，是团队工程能力的集中体现。第5章为测试与验收章节，通过Newman接口自动化测试与Playwright端到端运行结果对系统进行全面验证，与第4章形成建设与验证的闭环。第6章为总结与展望章节，凝练项目成果，分析当前技术局限，并从容器化部署、分布式事务与安全加固等方面提出后续改进方向。
